\documentclass[11pt]{article}

\usepackage[T1]{fontenc}
\usepackage{lmodern}
\usepackage[margin=1in]{geometry}
\usepackage{amsmath,amssymb,amsthm}
\usepackage{booktabs}
\usepackage{tabularx}
\usepackage{array}
\usepackage{microtype}
\usepackage{xcolor}
\usepackage{enumitem}
\usepackage{hyperref}

\hypersetup{
  colorlinks=true,
  linkcolor=blue!55!black,
  citecolor=blue!55!black,
  urlcolor=blue!55!black,
  pdftitle={Statewise versus operator Bell causality in finite quantum sequential growth},
  pdfauthor={Kenichi Osaki}
}

\setlist{nosep}
\setlength{\parindent}{0pt}
\setlength{\parskip}{0.45em}
\newtheorem{theorem}{Theorem}
\newtheorem{proposition}{Proposition}
\newtheorem{lemma}{Lemma}
\theoremstyle{definition}
\newtheorem{definition}{Definition}
\theoremstyle{remark}
\newtheorem{remark}{Remark}
\newcommand{\QQ}{\mathbb{Q}}
\newcommand{\GL}{\mathrm{GL}}
\newcommand{\End}{\mathrm{End}}
\newcommand{\im}{\operatorname{im}}
\newcommand{\draftmarker}[1]{%
  \begin{quote}\small
  \fcolorbox{red!60!black}{red!6}{%
    \parbox{0.90\linewidth}{\textbf{DRAFT PROOF INSERT REQUIRED.} #1}}
  \end{quote}}

\title{\textbf{Statewise versus operator Bell causality in finite\\
quantum sequential growth: exact separation and recovery at dimension two}}
\author{Kenichi Osaki\\
\small Independent researcher\\
\small ORCID: \href{https://orcid.org/0009-0003-9256-7089}{0009-0003-9256-7089}}
\date{Paper I manuscript scaffold, v0.4.2 -- 5 August 2026}

\begin{document}
\maketitle

\begin{center}
\small\emph{Claim-locked manuscript scaffold.  The theorem statements below
are limited to the v0.4.2 Paper I claim ledger.  Every proof section that
has not yet been written for human review carries an explicit draft marker.}
\end{center}

\begin{abstract}
We assemble a bounded manuscript for the finite, nonsingular, occurrence-ON
quantum sequential-growth presentation with \(d=2\) and source stages
\(n\leq4\).  The manuscript records exactly five ledger-bound statements:
a frozen strong/strong commutativity baseline; an exact rational
fixed-vector/general-covariance and reachable-state/martingale witness of
noncommutativity; a source-native Eq.~(120) lemma; commutativity on a proper
reconstruction slice; and conditional statewise-to-operator recovery lemmas.
The purpose is to keep the distinction between statewise and operator
semantics auditable, not to classify all finite semantics.  In particular,
the full 955 and 721 profiles, a complete finite ON classification, and a
reachable-visible witness are not claimed.  The bounded U2 auxiliary-ideal
investigation is retained only as a reproducibility and resource-boundary
record, not as a unit-ideal theorem.
\end{abstract}

\section{Introduction and claim boundary}

Classical sequential growth motivates covariance and Bell-causality
conditions on transitions between causal sets \cite{RideoutSorkin2000}.
The operator-valued CPOBC setting considered here follows the finite
presentation and source relations documented in \cite{SrivastavaSurya2026},
with the project-specific semantic distinctions kept explicit.

This draft is intentionally narrower than a general representation theorem.
It is a paper skeleton for the five statements C1--C5 in
\texttt{results/v0.4.2\_paper1\_claim\_boundary.json}.  The terms
``finite'', ``nonsingular'', ``occurrence ON'', \(d=2\), and \(n\leq4\)
are part of every theorem scope unless a statement explicitly says otherwise.
Exact certificates use arithmetic over \(\QQ\); their declared
characteristic-zero interpretation is retained without adding a new
field-extension claim.

Two distinctions organize the paper.  First, a residual may be required to
vanish as an operator, or only after evaluation on a specified state.
Second, a source assignment need not lie in the image of the frozen rational
reconstruction map.  The latter distinction is why the reconstruction-slice
theorem below is a proper-slice statement rather than a theorem for the full
strong-GC/reachable-state-MSR profile.

\section{Relation to adjacent rigidity results}

Xu's 2026 theorem package is materially broader on finite-dimensional
self-adjoint nonsingular systems and on several triangular and extension
branches \cite{Xu2026Relational}.  It also identifies non-self-adjoint
nonsingular transitions and state-only covariance as open boundaries.  The
present manuscript therefore makes no broad rigidity or priority claim over
that territory.  Its distinct center is the exact statewise weak/weak
separation in Theorem~\ref{thm:c2}, together with the explicitly conditional
recovery and proper-slice statements below.

The comparison in this draft is bound to the versioned literature audit
through 1 August 2026.  A fresh delta search is required before manuscript
freeze; an unavailable companion manuscript mentioned by Xu is not treated
as inspected.

\begin{table}[ht]
\centering
\caption{Claim boundary relative to the closest adjacent rigidity result.}
\label{tab:xu-comparison}
\begin{tabularx}{\textwidth}{@{}p{0.20\textwidth}p{0.34\textwidth}X@{}}
\toprule
Result & Declared setting & Manuscript use and boundary\\
\midrule
Xu \cite{Xu2026Relational}
  & self-adjoint nonsingular systems and several triangular/extension branches
  & broader in those settings; no project priority claim\\
C1 baseline
  & frozen literal strong-operator, \(d=2\), \(n\leq4\), occurrence ON
  & finite comparison endpoint for \(Q_1,\ldots,Q_4\) only\\
C2 separation
  & fixed-vector GC and reachable-state MSR over \(\QQ\)
  & exact non-self-adjoint statewise separation; reachable rank one\\
\bottomrule
\end{tabularx}
\end{table}

\section{Finite setup and semantic definitions}

\subsection{Frozen finite domain}

Let \(Q_1,\ldots,Q_4\in\GL_2(\QQ)\) denote the declared stage generators;
the statements below do not add a conclusion about \(Q_5\).  The transition
source cutoff is \(n\leq4\), while endpoint-stage-five data and \(Q_5\) may
still occur where a source formula requests them.  Every commutator conclusion
in this manuscript concerns \(Q_1,\ldots,Q_4\), and the C1 certificate does
not use \(Q_5\).  The transition domain is the declared nonsingular domain,
and occurrence identification is the certified ON domain of the underlying
artifacts.  Literal and derived Eq.~(113) inventories, and printed-strict and
completed Eq.~(139) inventories, remain separately recorded rather than
being identified by this manuscript.

\begin{definition}[Statewise and strong residual conditions]
For a same-endpoint path-pair residual \(D_{\alpha,\beta}\), fixed-vector
general covariance means
\[
  D_{\alpha,\beta}\Omega=0
\]
on the declared initial vector \(\Omega\); strong general covariance means
\(D_{\alpha,\beta}=0\) as an operator.  For a source martingale residual
\(M_c\) and its declared state \(v_c\), reachable-state MSR means
\[
  M_cv_c=0,
\]
while strong MSR means \(M_c=0\).  These definitions describe only the
declared finite residual inventories.
\end{definition}

\begin{definition}[Residual evaluation]
Let \(\mathcal R\subseteq\End(\QQ^2)\) be a declared linear space of
admissible residual operators, and let \(V\subseteq\QQ^2\) be a collection
of probes applied to the \emph{same} residual.  Define
\[
  \operatorname{ev}_V:\mathcal R\longrightarrow(\QQ^2)^V,
  \qquad D\longmapsto(Dv)_{v\in V}.
\]
The word ``recovery'' below always refers to this fixed residual family and
these sourcewise probes; it does not pool states tested against different
residual operators.
\end{definition}

\begin{definition}[Reconstruction slice]
Let
\(P_{\mathrm{sGC+rMSR}}^{\mathrm{ns}}\) be the set of source assignments
satisfying source-native CPOBC, strong GC, reachable-state MSR, and all
declared nonsingularity predicates.  Let \(U\) be the open domain on which
the four \(Q\)- and twenty-two \(B\)-inverse bases of the frozen
165-occurrence rational reconstruction are nonsingular, and let
\[
  \Phi_U:U\longrightarrow\text{(source assignments)}
\]
denote that reconstruction.  The reconstruction slice is
\(P_{\mathrm{sGC+rMSR}}^{\mathrm{ns}}\cap\im(\Phi_U)\).
\end{definition}

\begin{definition}[Reachable-visible noncommutativity]
A noncommutative witness is called reachable-visible here only if the
declared reachable sector exposes the nonzero commutator under the stated
visibility criterion.  This definition fixes a boundary term; the present
paper does not assert that its weak/weak witness has this property.
\end{definition}

\section{Ledger-bound theorem statements}

\subsection{C1: frozen strong/strong baseline}

\begin{theorem}[C1: finite strong/strong comparison theorem]
\label{thm:c1}
In the frozen v0.3.7/v0.3.9 literal strong-GC/strong-MSR finite
\(d=2\), \(n\leq4\), occurrence-ON baseline, the complete S1/S2/S3 cover
reduces the noncommutative alternative to 21 exact \(\QQ\) chart
calculations and the structural S3 certificate.  These exclude every
noncommutative branch, so every solution on the recorded denominator-open
locus satisfies
\[
  [Q_i,Q_j]=0\qquad(1\leq i<j\leq4).
\]
The certificate does not use \(Q_5\).
\end{theorem}

\draftmarker{Insert a human-readable account of the frozen literal
strong/strong certificate and its 21-chart-plus-S3 coverage.  Do not replace
this finite comparison theorem by a broad rigidity, self-adjoint, or complete
classification statement.}

\subsection{C2: exact weak/weak separation}

\begin{theorem}[C2: exact rational weak/weak separation]
\label{thm:c2}
The frozen finite \(d=2\), \(n\leq4\) presentation admits an exact rational,
nonsingular noncommutative witness satisfying fixed-vector GC and
reachable-state MSR.  The witness is an off-reachable-sector witness: its
reachable span has rank one, and it is not asserted to be reachable-visible.
The same orbit-constant assignment is also a witness with occurrence
identification OFF because the ON assignment is an admissible special case;
this does not supply a profile-native 406-occurrence OFF compiler or an OFF
classification.
\end{theorem}

\draftmarker{Insert the independently checkable witness derivation and the
exact finite audit.  Retain the distinction between statewise vanishing and
operator residuals; do not promote this theorem to either one-sided profile.}

\subsection{C3: source-native Eq.~(120)}

\begin{lemma}[C3: source-native Eq.~(120) lemma]
\label{lem:c3}
For units in a unital associative algebra, six raw antichain CPOBC relations
and the declared source nonsingularity imply the three \(k=1\) identities
\[
  Q_nQ_1^{-1}Q_m=Q_mQ_1^{-1}Q_n
  \qquad(2\leq m<n\leq4).
\]
The derivation uses neither MSR, GC, Eq.~(108), Eq.~(112), nor B-reduction.
It is stated independently of \(d=2\); dimension two enters only in the
subsequent chart closure.
\end{lemma}

\draftmarker{Insert the six selected source records and the free-word
derivation, with each cancellation and inverse operation stated.  The source
relation itself must be cited as prior material; the manuscript contribution
is the finite provenance and its use in the certified assembly.}

\subsection{C4: proper reconstruction slice}

\begin{theorem}[C4: reconstruction-slice commutativity]
\label{thm:c4}
For every
\[
  x\in P_{\mathrm{sGC+rMSR}}^{\mathrm{ns}}\cap\im(\Phi_U)
\]
and every \(q\in U\) with \(\Phi_U(q)=x\), the source-to-direct inclusion,
localisation, and the independently bound source relations give
\[
  q\in D_{955}.
\]
Lemma~\ref{lem:c3} and the exact \(d=2\) chart certificates then imply
\[
  [Q_i(q),Q_j(q)]=0\qquad(1\leq i<j\leq4).
\]
Consequently, the established forward image inclusion is
\[
  P_{\mathrm{sGC+rMSR}}^{\mathrm{ns}}\cap\im(\Phi_U)
  \subseteq\Phi_U(D_{955}),
\]
where \(D_{955}\) is the declared direct relation locus.  No converse
inclusion and no conclusion on
\(P_{\mathrm{sGC+rMSR}}^{\mathrm{ns}}\setminus\im(\Phi_U)\) are claimed.
\end{theorem}

\draftmarker{Insert the quantified source-to-direct pullback, the
localisation argument, and the exact 12-S1, 9-S2, and S3 closure.  Keep the
proper-slice qualifier in the theorem title, abstract, and conclusion.}

\subsection{C5: conditional statewise-to-operator recovery}

\begin{theorem}[C5: residual-space recovery]
\label{thm:c5}
For a declared residual space \(\mathcal R\) and sourcewise probe family
\(V\), statewise equality promotes every \(D\in\mathcal R\) to the operator
identity \(D=0\) if and only if \(\operatorname{ev}_V\) is injective.
In dimension two, two independent probes applied to the same residual are a
sufficient exact recovery condition.  The corresponding non-scalar
\(2\times2\) centralizer endpoint is
\[
  \operatorname{Cent}(A)=\QQ[I,A]
\]
for a non-scalar \(A\in M_2(\QQ)\).
\end{theorem}

\draftmarker{Insert the residual-space rank calculation, same-residual
two-probe argument, cyclicity counterexample, and centralizer proof.  Do not
assert that ordinary single-\(\Omega\) CPOBC semantics supplies the required
additional probes.}

\section{Proof roadmap and dependency audit}

The proof architecture is deliberately modular.  Lemma~\ref{lem:c3} is
source-native and dimension-independent.  The reconstruction argument in
Theorem~\ref{thm:c4} then binds source data to the direct locus, including
its localisation predicates.  Only after that binding do the \(d=2\)
S1/S2/S3 chart certificates close the noncommutative branches.  The
strong/strong baseline, Theorem~\ref{thm:c1}, remains a comparison theorem,
while Theorem~\ref{thm:c2} demonstrates that the two statewise conditions do
not themselves force operator commutativity.

\begin{center}
\small
raw source CPOBC \(\longrightarrow\) Eq.~(120)
\(\longrightarrow\) source-to-direct inclusion and localisation
\(\longrightarrow\) commuting ratios
\(\longrightarrow\) S1/S2/S3 chart closure.
\end{center}

The principal certificate census for the reconstruction-slice path is
shown in Table~\ref{tab:spine}.  Counts are evidence ledgers, not a claim of
relation minimality.

\begin{table}[ht]
\centering
\caption{Certified proof-spine census used by the proper-slice statement.}
\label{tab:spine}
\begin{tabularx}{\textwidth}{@{}l r X@{}}
\toprule
Item & Count & Role in the manuscript\\
\midrule
source CPOBC records & \(783=83+700\) & source-to-direct CPOBC pullback\\
strong-GC generating basis & \(320=65+255\) & direct local-operator-GC pullback\\
same-endpoint path pairs & 1,529 & coverage for the declared GC basis\\
reconstructed transitions & 165 & nonsingularity and occurrence binding\\
rational inverse bases & \(4\ Q+22\ B\) & localisation origin\\
raw Eq.~(120) records / targets & \(6/3\) & source-native lemma\\
dimension-two cover & \(12\ \mathrm{S1}+9\ \mathrm{S2}+\mathrm{S3}\) & chart closure\\
exact \(\QQ\) certificates used & 21 & proper-slice closure only\\
\bottomrule
\end{tabularx}
\end{table}

\draftmarker{This section is a dependency roadmap, not a substitute for the
complete proofs.  Before release, add the selected record IDs, quantified
map domains, human-readable chart-cover proof, and certificate-to-claim
bindings.}

\section{Witness audit for C2}

The C2 witness is exact over \(\QQ\).  With
\(\Omega=(1,0)^\mathsf{T}\), the recorded transition family has the form
\[
 A_e=
 \begin{pmatrix}2^{w-m}/2^n & b_e\\0&1\end{pmatrix},
 \qquad
 b_e=\begin{cases}4/2^n,&\text{for a gregarious transition},\\
                    0,&\text{otherwise},\end{cases}
\]
where \(w\) and \(m\) are the declared precursor statistics.  For the
antichain gregarious generators,
\[
 Q_n=\begin{pmatrix}2^{-n}&4\,2^{-n}\\0&1\end{pmatrix}.
\]
Thus the recorded witness has nonzero commutators among
\(Q_1,\ldots,Q_4\), while every reachable state is a nonzero rational
multiple of \(\Omega\).  This last fact is the reason its noncommutativity
is not called reachable-visible.

\begin{table}[ht]
\centering
\caption{Exact finite checks recorded for the C2 witness.}
\label{tab:witness}
\begin{tabular}{lr}
\toprule
Audit family & Recorded checks\\
\midrule
frozen CPOBC word equations & 783\\
same-endpoint fixed-vector GC pairs & 1,529\\
reachable-state MSR sources & 24\\
Eq.~(113), derived and literal branches & \(25+25\)\\
Eq.~(139), printed-strict and completed domains & \(4+10\)\\
transition occurrence determinants & 165\\
nonzero commutators among \(Q_1,\ldots,Q_4\) & 6\\
\bottomrule
\end{tabular}
\end{table}

\draftmarker{Before submission, replace this audit summary by the complete
witness table, direct residual ledger, and independent-oracle comparison.
Do not infer a complete weak/weak classification from this witness.}

\section{Recovery conditions and their boundary}

Theorem~\ref{thm:c5} is deliberately conditional.  One probe cannot
separate all of \(M_2(\QQ)\): evaluation at \(e_1\), for example, has a
nontrivial kernel on the full matrix algebra.  Conversely, the basis probes
\(e_1,e_2\) recover all of \(M_2(\QQ)\), and a single probe can separate a
smaller declared residual space.  The correct condition is therefore
injectivity on \(\mathcal R\), not a generic slogan about cyclicity or global
reachability.

The sourcewise qualifier cannot be omitted.  The exact counterexample
\[
  v_1=e_1,\quad D_1=E_{12},\qquad
  v_2=e_2,\quad D_2=E_{21}
\]
has \(\operatorname{span}\{v_1,v_2\}=\QQ^2\) and
\(D_1v_1=D_2v_2=0\), although both residuals are nonzero.  Vectors tested
against distinct source residuals therefore cannot be pooled to recover
either residual as an operator identity.

\draftmarker{Add the exact evaluation-rank ledger and make the bridge to the
strong/strong baseline only under separately verified injectivity for each
declared GC and MSR residual family.}

\section{Explicit nonclaims}

The following boundaries are part of the manuscript, rather than future
work inferred from its results:
\begin{enumerate}[label=N\arabic*.]
  \item No general strong-GC/reachable-state-MSR theorem for the full
        955 profile is claimed.
  \item No general fixed-vector-GC/strong-MSR theorem for the 721 profile
        is claimed.
  \item No complete finite occurrence-ON semantics lattice is claimed.
  \item The current weak/weak witness is not claimed to be
        reachable-visible.
  \item The bounded GF determinantal U2 scout is not a \(\QQ\) theorem that
        the full auxiliary ideal is unit or non-unit.
  \item A general or profile-native occurrence-OFF compiler and
        classification, relation minimality, \(n\geq5\), \(d\geq3\), singular
        transitions, and infinite-system claims are out of scope.  The C2
        special-case OFF witness does not remove this boundary.
\end{enumerate}

The global SR2-V status remains
\texttt{SEARCH\_OPEN\_NO\_TERMINAL}.  The U2 subroute status
\texttt{SOFT\_RESOURCE\_LIMIT\_NONTERMINAL} records a bounded computation and
verified process cleanup; it is neither an obstruction theorem nor a terminal
verdict.

\section{Reproducibility appendix}

Table~\ref{tab:artifacts} records the claim-bearing machine artifacts.  The
manuscript must preserve their raw hashes and semantic digests rather than
substituting prose summaries for them.

\begin{table}[ht]
\centering
\caption{Paper I claim ledger and principal evidence artifacts.}
\label{tab:artifacts}
\small
\begin{tabularx}{\textwidth}{@{}p{0.14\textwidth}p{0.32\textwidth}X@{}}
\toprule
Claim & Artifact & Recorded binding\\
\midrule
C1 & \texttt{results/v0.3.7\_release\_manifest.json};
     \texttt{results/v0.3.9\_release\_manifest.json}
   & frozen public baseline; v0.3.7 semantic digest
     \texttt{b97b83\ldots bc58d}; v0.3.9 semantic digest
     \texttt{eb9eb8\ldots 48e47}\\
C2 & \texttt{results/v0.4\_weak\_d2\_classification.json};
     \texttt{results/v0.4.2\_sr2v\_baseline\_observability.json}
   & witness raw SHA-256 \texttt{18e71f\ldots 2fe65}; observability raw
     SHA-256 \texttt{4f5780\ldots 10de}; semantic digest
     \texttt{73508f\ldots c62}\\
C3 & \texttt{results/v0.4.2\_eq120\_source\_provenance.json}
   & raw SHA-256 \texttt{864536\ldots fed4}; semantic digest
     \texttt{77d11e\ldots d6ca}\\
C4 & \texttt{results/v0.4.1\_partial\_slice\_audit.json}
   & raw SHA-256 \texttt{a37bc2\ldots 089b5}; semantic digest
     \texttt{2d65dd\ldots 99839}\\
C5 & \texttt{results/v0.4.2\_sr3b\_m\_semantic\_recovery.json}
   & raw SHA-256 \texttt{164dba\ldots 235a8}; semantic digest
     \texttt{a74297\ldots 5d735}\\
\bottomrule
\end{tabularx}
\end{table}

The current reproduction-facing commands are:
\begin{verbatim}
uv run pytest -q tests/final_theory/test_eq120_source_provenance_v042.py
uv run python scripts/reproduce_v041_partial_slice_audit.py --write-audit
uv run pytest -q tests/final_theory/test_v041_partial_slice_audit.py
python -m universe_lab.final_theory.semantic_recovery_v042
pytest -q tests/final_theory/test_semantic_recovery_v042.py
\end{verbatim}

The resource-boundary record for the U2 subroute is
\texttt{reports/v0.4.2\_sr2v\_auxiliary\_ideal\_determinantal\_pilot\_2026-08-05.md}.
It has no theorem role beyond reproducibility and limitation reporting.

\draftmarker{Before a release candidate, verify these commands in a clean
environment, expand abbreviated digests to their full 64 hexadecimal
characters in the release manifest, and add the versioned software and
resource ledger.}

\section*{Manuscript status}

This file is a scoped assembly draft.  It is not a submission, deposit, or
claim that the required human-readable proofs, independent audits, and
release checks have been completed.  The owner and manuscript-review gates
remain pending.

\bibliographystyle{plain}
\bibliography{../v0.3.9_d2_commutativity_short_report/references,references}

\end{document}
